
\documentclass{article}
\usepackage[utf8]{inputenc}
\usepackage[T1]{fontenc}

\usepackage{amsmath,amssymb}

\title{Définition d'une Variété (Manifold)}


\begin{document}
\maketitle

Une \textbf{variété différentiable} de dimension $n$ est un espace topologique $M$ tel que pour tout $p \in M$, il existe un ouvert $U \subset M$ contenant $p$ et un homéomorphisme $\phi: U \to V \subset \mathbb{R}^n$ (où $V$ est ouvert). Le couple $(U,\phi)$ est une \textbf{carte de coordonnées}.

\textbf{Idées clés:}
\begin{itemize}
    \item \textit{Localité:} localement $\cong \mathbb{R}^n$, mais topologie globale peut être non-triviale (sphère, tore)
    \item \textit{Dimension:} $n$ coordonnées nécessaires localement
    \item \textit{Cartes:} plusieurs cartes couvrent la variété entière
\end{itemize}

\textbf{Exemples:} $S^1$ (dim 1), $S^2$ (dim 2), spacetime GR (dim 4, métrique courbe).

\textbf{Conclusion (Carroll):} les variétés sont \textit{la scène} de la relativité générale — espaces courbes avec calcul différentiel local possible.

\end{document}
